\begin{lemma}[Bounded sequences in $L^2(\gamma_n)$ have weakly convergent subsequences]
Let $(v_k)_{k \geq 0}$ be a sequence in $L^2(\gamma_n)$ and let $R \in \mathbb{R}$
satisfy $\|v_k\|_{L^2(\gamma_n)} \leq R$ for every $k$.  Then there are indices
$k_0 < k_1 < k_2 < \cdots$ and an element $u \in L^2(\gamma_n)$ with
$\|u\|_{L^2(\gamma_n)} \leq R$ such that
\[
  \langle v_{k_j},\, \varphi\rangle_{L^2(\gamma_n)}
    \longrightarrow \langle u,\, \varphi\rangle_{L^2(\gamma_n)}
  \qquad\text{as } j \to \infty, \text{ for every } \varphi \in L^2(\gamma_n).
\]
This is weak sequential compactness of bounded sets in the single Hilbert space
$L^2(\gamma_n)$.  No weak topology on $\mathcal{V} = \operatorname{Dom}(L^{1/4})$
is used, no reflexivity of $\mathcal{V}$ is invoked, and no compactness of any
embedding and no dependence on the dimension $n$ is asserted.
\end{lemma}
\begin{proof}
Write $H = L^2(\gamma_n)$, a real Hilbert space, and $H^*$ for its dual.  The
proof has two steps: $H$ is separable, and the weak-$*$ statement for $H^*$
recorded in the shared context transports to $H$ through the Riesz
representation.  If $R < 0$ the hypothesis is vacuous, since
$\|v_0\| \geq 0 > R$ is impossible; so $R \geq 0$ throughout.

\emph{Step 1: the index set $\mathbb{N}_0^n$ is countable.}
$\mathbb{N}_0^n$ is the set of functions $\{1,\dots,n\} \to \mathbb{N}_0$, i.e.\ a
product of $n$ copies of the countable set $\mathbb{N}_0$.  A product of finitely
many countable sets is countable, by induction on the number of factors: the
empty product is a one-point set, and if $A$ is countable and $B$ is countable
then $A \times B$ is countable because $\mathbb{N}_0 \times \mathbb{N}_0$ is
countable (e.g.\ by the injection $(i,j) \mapsto 2^i 3^j$) and a set injecting
into a countable set is countable.  Hence $\mathbb{N}_0^n$ is countable, and we
may fix an enumeration of it.

\emph{Step 2: $H$ is separable.}
Let
\[
  D = \Bigl\{\textstyle\sum_{\alpha \in F} q_\alpha\, h_\alpha
      \;:\; F \subseteq \mathbb{N}_0^n \text{ finite},\;
            q_\alpha \in \mathbb{Q} \text{ for } \alpha \in F \Bigr\}
    \subseteq H ,
\]
with $h_\alpha$ the Hermite functions of \noderef{HermiteL2}.

$D$ is countable.  By Step 1 the set of finite subsets of $\mathbb{N}_0^n$ is
countable (the finite subsets of a countable set form a countable union, over
$m \in \mathbb{N}_0$, of the images of the countable sets
$(\mathbb{N}_0^n)^m$ under the map sending an $m$-tuple to the set of its
entries), and for each fixed finite $F$ the set of families
$(q_\alpha)_{\alpha \in F} \in \mathbb{Q}^F$ is countable, $\mathbb{Q}$ being
countable and $F$ finite.  So $D$ is a countable union of countable sets, hence
countable.

$D$ is dense in $H$.  Let $v \in H$ and $\varepsilon > 0$.  By
\noderef{HermiteL2Complete} the closed linear span of
$\{h_\alpha : \alpha \in \mathbb{N}_0^n\}$ is all of $H$, so $v$ lies in the
closure of the set of finite real linear combinations of the $h_\alpha$; hence
there are a finite $F \subseteq \mathbb{N}_0^n$ and reals $(c_\alpha)_{\alpha\in F}$
with
\[
  \Bigl\| v - \sum_{\alpha \in F} c_\alpha h_\alpha \Bigr\| < \frac{\varepsilon}{2}.
\]
If $F = \emptyset$ the element $0 \in D$ already satisfies
$\|v - 0\| < \varepsilon/2$, so assume $F \neq \emptyset$ and let $N = |F| \geq 1$.
Since $\mathbb{Q}$ is dense in $\mathbb{R}$, choose for each $\alpha \in F$ a
rational $q_\alpha$ with $|c_\alpha - q_\alpha| < \varepsilon/(2N)$, and put
$w = \sum_{\alpha\in F} q_\alpha h_\alpha \in D$.  By
\noderef{HermiteL2Orthonormal} each $h_\alpha$ is a unit vector,
$\|h_\alpha\| = 1$, so by the triangle inequality in $H$
\[
  \Bigl\| \sum_{\alpha\in F} c_\alpha h_\alpha - w \Bigr\|
    = \Bigl\| \sum_{\alpha\in F} (c_\alpha - q_\alpha) h_\alpha \Bigr\|
    \leq \sum_{\alpha\in F} |c_\alpha - q_\alpha|\,\|h_\alpha\|
    < N \cdot \frac{\varepsilon}{2N} = \frac{\varepsilon}{2}.
\]
Adding the two displays, $\|v - w\| < \varepsilon$ with $w \in D$.  As $v$ and
$\varepsilon$ were arbitrary, $D$ is a countable dense subset of $H$, i.e.\ $H$
is separable.

\emph{Step 3: the sequence of associated functionals.}
For each $k$ define $\Lambda_k : H \to \mathbb{R}$ by
$\Lambda_k(w) = \langle v_k, w\rangle$.  Each $\Lambda_k$ is linear in $w$ by
linearity of the inner product in its second argument, and by Cauchy--Schwarz,
which is shared context, $|\Lambda_k(w)| \leq \|v_k\|\,\|w\|$ for every $w$, so
$\Lambda_k$ is continuous with $\|\Lambda_k\|_{H^*} \leq \|v_k\|$.  Thus
$\Lambda_k \in H^*$, and applying the Riesz representation item of the shared
context to $\Lambda_k$ produces a \emph{unique} $y \in H$ with
$\Lambda_k = \langle y, \cdot\rangle$ and $\|\Lambda_k\|_{H^*} = \|y\|$; since
$v_k$ is such an element by the definition of $\Lambda_k$, uniqueness gives
$y = v_k$ and therefore
\begin{equation*}
  \|\Lambda_k\|_{H^*} = \|v_k\| \leq R \qquad\text{for every } k. \tag{1}
\end{equation*}

\emph{Step 4: extraction, and transport back to $H$.}
By Step 2 the normed space $H$ is separable, and by $(1)$ the sequence
$(\Lambda_k)$ lies in the ball of radius $R$ of $H^*$.  The weak-$*$ sequential
compactness item of the shared context, applied to $E = H$, this $R \geq 0$ and
this sequence, therefore yields indices $k_0 < k_1 < k_2 < \cdots$ and
$\Lambda \in H^*$ with $\|\Lambda\|_{H^*} \leq R$ such that
\begin{equation*}
  \Lambda_{k_j}(w) \longrightarrow \Lambda(w)
  \qquad\text{for every } w \in H. \tag{2}
\end{equation*}
Applying the Riesz representation item once more, now to $\Lambda$, gives
$u \in H$ with $\Lambda(w) = \langle u, w\rangle$ for every $w \in H$ and
$\|u\| = \|\Lambda\|_{H^*} \leq R$.

Finally, let $\varphi \in H$ be arbitrary.  By the definition of $\Lambda_{k_j}$,
by $(2)$ with $w = \varphi$, and by the representation of $\Lambda$,
\[
  \langle v_{k_j},\, \varphi\rangle
    = \Lambda_{k_j}(\varphi)
    \longrightarrow \Lambda(\varphi)
    = \langle u,\, \varphi\rangle .
\]
So the indices $k_0 < k_1 < \cdots$ and the element $u$, with $\|u\| \leq R$,
are as asserted.
\end{proof}
